% !PRINTED 191
\weberpage{191}

Wenn wir diese Formeln durch $\vartheta_{01}(u)^{n^2}$ dividieren,
so lassen sich die rechten Seiten als ganze rationale Funktionen
von den elliptischen Funktionen $\mathrm{sn}\,v$, $\mathrm{cn}\,v$,
$\mathrm{dn}\,v$ darstellen (§ 42). Wenn wir also

\weq{3}{v = \pi\vartheta_{00}^2u, \quad \mathrm{sn}\,v=x, \quad \mathrm{cn}\,v=y, \quad \mathrm{dn}\,v=z}

setzen, so können wir die Formeln (1), (2), indem wir einen
konstanten Faktor passend bestimmen, so schreiben:

I. bei geradem $n$:

\[
\begin{array}{ll}
\dfrac{\vartheta_{01}^{n^2}\vartheta_{00}}{\vartheta_{10}}\,\dfrac{\vartheta_{11}(nu)}{\vartheta_{01}(u)^{n^2}} = xyz\,A(x^2), & A(0)=n,\\[3mm]
\dfrac{\vartheta_{01}^{n^2}}{\vartheta_{10}}\,\dfrac{\vartheta_{10}(nu)}{\vartheta_{01}(u)^{n^2}} = B(x^2), & B(0)=1,\\[3mm]
\dfrac{\vartheta_{01}^{n^2}}{\vartheta_{00}}\,\dfrac{\vartheta_{00}(nu)}{\vartheta_{01}(u)^{n^2}} = C(x^2), & C(0)=1,\\[3mm]
\vartheta_{01}^{n^2-1}\,\dfrac{\vartheta_{01}(nu)}{\vartheta_{01}(u)^{n^2}} = D(x^2), & D(0)=1.
\end{array}
\]

II. bei ungeradem $n$:

\[
\begin{array}{ll}
\dfrac{\vartheta_{00}\vartheta_{01}^{n^2-1}}{\vartheta_{10}}\,\dfrac{\vartheta_{11}(nu)}{\vartheta_{01}(u)^{n^2}} = x\,A(x^2), & A(0)=n,\\[3mm]
\dfrac{\vartheta_{01}^{n^2}}{\vartheta_{10}}\,\dfrac{\vartheta_{10}(nu)}{\vartheta_{01}(u)^{n^2}} = y\,B(x^2), & B(0)=1,\\[3mm]
\dfrac{\vartheta_{01}^{n^2}}{\vartheta_{00}}\,\dfrac{\vartheta_{00}(nu)}{\vartheta_{01}(u)^{n^2}} = z\,C(x^2), & C(0)=1,\\[3mm]
\vartheta_{01}^{n^2-1}\,\dfrac{\vartheta_{01}(nu)}{\vartheta_{01}(u)^{n^2}} = D(x^2), & D(0)=1,
\end{array}
\]

worin $A$, $B$, $C$, $D$ ganze rationale Funktionen von $x^2$ sind,
deren Grade sich aus den Formeln (1), (2) folgendermaßen ergeben:

% [sic] In the odd-n row the source prints (n^2-2)/2 for the degree of D, which is nonintegral for odd n; the defining formula implies (n^2-1)/2.
\[
\begin{array}{c@{\quad}cccc@{\qquad}l}
 & A(x^2), & B(x^2), & C(x^2), & D(x^2), & \\[3mm]
\text{Grad:} & \dfrac{n^2}{2}-2, & \dfrac{n^2}{2}, & \dfrac{n^2}{2}, & \dfrac{n^2}{2} & n \text{ gerade,}\\[3mm]
 & \dfrac{n^2-1}{2}, & \dfrac{n^2-1}{2}, & \dfrac{n^2-1}{2}, & \dfrac{n^2-2}{2} & n \text{ ungerade.}
\end{array}
\]

Aus den Definitionen I., II. erhält man unmittelbar die Werte von
$A(0)$, $B(0)$, $C(0)$, $D(0)$, d.\,h. die von $n$ unabhängigen
% [sic] The source calls these terms independent of n although the displayed value A(0)=n depends on n.
Glieder, wie sie beigefügt sind.

Mannigfache Beziehungen zwischen diesen Funktionen ergeben sich noch
auf folgendem Wege:
%CONT>>
